\exercisesection{\S~3}

\begin{enumerate}
\def\labelenumi{\arabic{enumi})}
\item
  Let g be a Lie algebra, \(g'\) a Cartan subalgebra of g. Then the conditions of Prop. 3 are satisfied. But an element of \(g'\) , even if it is regular in \(g'\) , is not necessarily regular in g.
\item
  Let g be a real Lie algebra of dimension n, U (resp. H) the set of regular elements (resp. of Cartan subalgebras) of g, and \(\text{Int}(\mathfrak{g})\) the group of inner automorphisms of g (Chap. III, §6, no. 2, Def. 2).
\end{enumerate}

\begin{enumerate}
\def\labelenumi{\alph{enumi})}
\item
  Show that, if \(x\) and \(y\) belong to the same connected component of \(\mathrm{U}\) , \(\mathfrak{g}^0(x)\) and \(\mathfrak{g}^0(y)\) are conjugate under \(\operatorname{Int}(\mathfrak{g})\) .
\item
  Show that the number of connected components of U is finite, and that this number is bounded by a constant \(c(n)\) depending only on n (apply Exerc. 2 of App. II).
\item
  Deduce that the number of orbits of \(\operatorname{Int}(\mathfrak{g})\) on \(\mathrm{H}\) is \(\leq c(n)\) .
\end{enumerate}

\begin{enumerate}
\def\labelenumi{\arabic{enumi})}
\setcounter{enumi}{2}
\tightlist
\item
  Let \(\mathfrak{g}\) be a real Lie algebra, \(\mathfrak{r}\) the radical of \(\mathfrak{g}\) , \(\mathfrak{h}\) and \(\mathfrak{h}'\) Cartan subalgebras of \(\mathfrak{g}\) , \(\varphi\) the canonical homomorphism from \(\mathfrak{g}\) to \(\mathfrak{g} / \mathfrak{r}\) . The following conditions are equivalent:
\end{enumerate}

\begin{enumerate}
\def\labelenumi{(\roman{enumi})}
\item
  \(\mathfrak{h}\) and \(\mathfrak{h}'\) are conjugate under \(\operatorname{Int}(\mathfrak{g})\) ;
\item
  \(\varphi(\mathfrak{h})\) and \(\varphi(\mathfrak{h}')\) are conjugate under \(\operatorname{Int}(\mathfrak{g}/\mathfrak{r})\) . (Imitate the proof of Prop. 5.)
\end{enumerate}

\begin{enumerate}
\def\labelenumi{\arabic{enumi})}
\setcounter{enumi}{3}
\item
  Let \(\mathfrak{g} = \mathfrak{sl}(2, \mathbf{R})\) , \(x = \begin{pmatrix} 1 & 0 \\ 0 & -1 \end{pmatrix}\) , \(y = \begin{pmatrix} 0 & -1 \\ 1 & 0 \end{pmatrix}\) . Show that \(\mathbf{R}x\) and \(\mathbf{R}y\) are Cartan subalgebras of \(\mathfrak{g}\) not conjugate under \(\operatorname{Aut}(\mathfrak{g})\) .
\item
  \begin{enumerate}
  \def\labelenumii{\alph{enumii})}
  \tightlist
  \item
    Show that there exists a Lie algebra \(\mathfrak{g}\) over \(k\) with basis \((x,y,z,t)\) such that
  \end{enumerate}
\end{enumerate}

\[
[ x, y ] = z, [ x, t ] = t, [ y, t ] = 0, [ \mathfrak {g}, z ] = 0.
\]

Show that g is solvable and that \(k = kx + ky + kz\) is a subalgebra of g.

\begin{enumerate}
\def\labelenumi{\alph{enumi})}
\setcounter{enumi}{1}
\item
  Show that the elementary automorphisms of g are the maps of the form \(1 + \lambda \operatorname{ad}_{\mathfrak{g}} y + \mu \operatorname{ad}_{\mathfrak{g}} t\) where \(\lambda, \mu \in k\) .
\item
  Show that \(1 + ad_{k}x\) is an elementary automorphism of k which does not extend to an elementary automorphism of g.
\end{enumerate}

\(d\) ) Let \(\mathfrak{s}\) be a semi-simple subalgebra of a Lie algebra \(\mathfrak{a}\) . Show that every elementary automorphism of \(\mathfrak{s}\) extends to an elementary automorphism of \(\mathfrak{a}\) .

\begin{enumerate}
\def\labelenumi{\arabic{enumi})}
\setcounter{enumi}{5}
\item
  Every element of a reductive Lie algebra \(\mathfrak{g}\) is contained in a commutative subalgebra of dimension \(\mathrm{rk}(\mathfrak{g})\) .
\item
  Let \(\mathfrak{g}\) be a Lie algebra and \(\mathfrak{g}'\) a subalgebra of \(\mathfrak{g}\) reductive in \(\mathfrak{g}\) . Let \(\mathfrak{a}\) be a Cartan subalgebra of \(\mathfrak{g}'\) . Show that there exists a Cartan subalgebra of \(\mathfrak{g}\) which contains \(\mathfrak{a}\) (use Prop. 10 of §2). Deduce that \(\mathrm{rk}(\mathfrak{g}') \leq \mathrm{rk}(\mathfrak{g})\) and that equality holds if and only if \(\mathfrak{g}'\) has properties (i), (ii), (iii) of Prop. 3.
\item
  Let \(\mathfrak{g}\) be a Lie algebra, \(\mathfrak{a}\) an ideal of \(\mathfrak{g}\) , \(\mathfrak{h}\) a Cartan subalgebra of \(\mathfrak{g}\) and \(\mathcal{C}_{\infty}\mathfrak{g}\) the union of the ascending central series of \(\mathfrak{g}\) . Show that \(\mathfrak{a} \subset \mathfrak{h}\) implies \(\mathfrak{a} \subset \mathcal{C}_{\infty}\mathfrak{g}\) (in other words \(\mathcal{C}_{\infty}\mathfrak{g}\) is the largest ideal of \(\mathfrak{g}\) contained in \(\mathfrak{h}\) ). (Reduce to the case where \(k\) is algebraically closed and remark that \(\mathfrak{a}\) is stable under every elementary automorphism of \(\mathfrak{g}\) ; the relation \(\mathfrak{a} \subset \mathfrak{h}\) then implies that \(\mathfrak{a}\) is contained in every Cartan subalgebra of \(\mathfrak{g}\) ; conclude by means of Exerc. 15 of §2.)
\end{enumerate}

¶ 9) Let g be a Lie algebra, h a Cartan subalgebra of g and x an element of h. Let \(g = h \oplus g^{+}\) be the Fitting decomposition (§1, no. 1) of g with respect to h.

\begin{enumerate}
\def\labelenumi{\alph{enumi})}
\item
  Let \(\mathfrak{n}\) be the largest semi-simple \(\mathfrak{h}\) -submodule contained in \(\mathfrak{g}^0 (x)\cap \mathfrak{g}^+\) . Show that \(\mathfrak{n} = 0\) if and only if \(\mathfrak{g}^0 (x) = \mathfrak{h}\) , i.e.~\(x\) is regular in \(\mathfrak{g}\) .
\item
  Show that \(\mathfrak{h} \oplus \mathfrak{n}\) is a subalgebra of \(\mathfrak{g}\) . If \(\mathfrak{h}'\) is the intersection of \(\mathfrak{h}\) with the commutant of \(\mathfrak{n}\) , show that \(\mathfrak{h}'\) is an ideal of \(\mathfrak{h} \oplus \mathfrak{n}\) which contains \(\mathcal{D}\mathfrak{h}\) and \(x\) . Conclude (Exerc. 8) that \(\mathfrak{h}' \subset \mathcal{C}_{\infty}(\mathfrak{h} \oplus \mathfrak{n})\) , and hence that \(x \in \mathcal{C}_{\infty}(\mathfrak{h} \oplus \mathfrak{n})\) and that \(x\) belongs to every Cartan subalgebra of \(\mathfrak{h} \oplus \mathfrak{n}\) .
\item
  If \(n \neq 0\) , \(h \oplus n\) is not nilpotent and has infinitely-many Cartan subalgebras (§2, Exerc. 10). Conclude that x belongs to infinitely-many Cartan subalgebras of g.
\end{enumerate}

\(d\) ) An element of \(\mathfrak{g}\) is regular if and only if it belongs to a unique Cartan subalgebra.

¶ 10) Let g be a Lie algebra, r its radical, n its largest nilpotent ideal and a one of its Levi subalgebras.

\begin{enumerate}
\def\labelenumi{\alph{enumi})}
\item
  Put \(\mathfrak{g}' = \mathfrak{n} + \mathcal{D}\mathfrak{g}\) . Show that \(\mathfrak{g}' = \mathfrak{n} \oplus \mathfrak{a}\) . (Use the fact that \([\mathfrak{g}, \mathfrak{r}]\) is contained in \(\mathfrak{n}\) .) If \(\mathfrak{g} \neq 0\) , then \(\mathfrak{g}' \neq 0\) .
\item
  Assume that k is algebraically closed. Let \((\mathrm{V}_{i})_{i\in\mathrm{I}}\) be the quotients of a Jordan-Hölder sequence of the g-module g (with the adjoint representation). If \(x\in r\) , show that \(x_{V_{i}}\) is a homothety and that \(x_{V_{i}}=0\) for all i if and only if x belongs to n.~Deduce that an element \(y\in g\) belongs to \(g'\) if and only if \(\operatorname{Tr}(y_{\mathrm{V}_{i}})=0\) for all \(i\in I\) .
\item
  Denote by N the vector subspace of g generated by the elements x such that ad x is nilpotent. Show that N is a subalgebra of g (use the fact that N is stable under \(\operatorname{Aut}_{e}(\mathfrak{g})\) ). Show, by using b), that \(N \subset g'\) .
\end{enumerate}

\(d\) ) Let \(\mathfrak{h}\) be a Cartan subalgebra of \(\mathfrak{g}\) . Assume that there exists a subset \(R\) of \(\mathfrak{h}^*\) such that

\[
\mathfrak {g} = \mathfrak {h} \oplus \bigoplus_ {\alpha \in R} \mathfrak {g} ^ {\alpha} (\mathfrak {h}),
\]

an assumption which is satisfied, in particular, if k is algebraically closed. Show that N then contains the \(\mathfrak{g}^{\alpha}(\mathfrak{h})\) , Dh and n; deduce that N contains \(g'\) , so \(N = g'\) .

\begin{enumerate}
\def\labelenumi{\alph{enumi})}
\setcounter{enumi}{4}
\tightlist
\item
  If \(k\) is algebraically closed and \(\mathfrak{g} \neq 0\) , \(\mathfrak{g}\) contains an element \(x \neq 0\) such that \(\operatorname{ad} x\) is nilpotent. (Indeed, we then have \(\mathfrak{g}' \neq 0\) .)
\end{enumerate}

¶ 11) Let g be a Lie algebra, r its radical.

\begin{enumerate}
\def\labelenumi{\alph{enumi})}
\item
  Let \(\mathfrak{s}\) be a Levi subalgebra of \(\mathfrak{g}\) , \(\mathfrak{k}\) a Cartan subalgebra of \(\mathfrak{s}\) . Show that \(\mathfrak{k}\) is contained in a Cartan subalgebra \(\mathfrak{h}\) of \(\mathfrak{g}\) which is the sum of \(\mathfrak{k}\) and a subalgebra of \(\mathfrak{r}\) . (Use §2, Th. 2, Prop. 10 and Cor. 2 of Th. 1.)
\item
  Let \(\mathfrak{h}'\) be a Cartan subalgebra of \(\mathfrak{g}\) . Show that there exists a Levi subalgebra \(\mathfrak{s}'\) of \(\mathfrak{g}\) such that \(\mathfrak{h}'\) is the sum of a Cartan subalgebra of \(\mathfrak{s}'\) and a subalgebra of \(\mathfrak{r}\) . (The subalgebras \(\mathfrak{s}, \mathfrak{k}, \mathfrak{h}\) in \(a\) ) can be chosen so that \(\mathfrak{h} + \mathfrak{r} = \mathfrak{h}' + \mathfrak{r}\) . Put \(\mathfrak{a} = \mathfrak{h} + \mathfrak{r}\) , which is solvable. By Th. 3, there exists \(x \in \mathcal{C}^{\infty}(\mathfrak{a})\) such that \(e^{\mathrm{ad}_{\mathfrak{a}}x}\mathfrak{h} = \mathfrak{h}'\) . Then \(e^{\mathrm{ad}_{\mathfrak{g}}x}\) is a special automorphism of \(\mathfrak{g}\) which transforms \(\mathfrak{s}\) into the required Levi subalgebra.)
\item
  Let \(\mathfrak{s}\) be a Levi subalgebra of \(\mathfrak{g}\) . Let \(\mathfrak{h}\) be a Cartan subalgebra of \(\mathfrak{g}\) which is the sum of a Cartan subalgebra \(\mathfrak{k}\) of \(\mathfrak{s}\) and a subalgebra \(\mathfrak{l}\) of \(\mathfrak{r}\) . Let \(\mathfrak{c}\) be the commutant of \(\mathfrak{k}\) in \(\mathfrak{r}\) . Show that \(\mathfrak{l}\) is a Cartan subalgebra of \(\mathfrak{c}\) . (For \(x \in \mathfrak{k}\) , \(\operatorname{ad}_{\mathfrak{g}} x\) is semi-simple, but \(\operatorname{ad}_{\mathfrak{h}} x\) is nilpotent, so \([\mathfrak{k}, \mathfrak{h}] = 0\) . If \(y \in \mathfrak{c}\) is such that \([y, \mathfrak{l}] \subset \mathfrak{l}\) , then \([y, \mathfrak{h}] \subset \mathfrak{h}\) so \(y \in \mathfrak{h} \cap \mathfrak{r} = \mathfrak{l}\) .)
\item
  Let \(\mathfrak{s}\) be a Levi subalgebra of \(\mathfrak{g}\) , \(\mathfrak{k}\) a Cartan subalgebra of \(\mathfrak{s}\) , \(\mathfrak{c}\) the commutant of \(\mathfrak{k}\) in \(\mathfrak{r}\) , \(\mathfrak{l}\) a Cartan subalgebra of \(\mathfrak{c}\) . Then \(\mathfrak{h} = \mathfrak{k} + \mathfrak{l}\) is a Cartan subalgebra of \(\mathfrak{g}\) . (Let \(x = y + z\) ( \(y \in \mathfrak{s}, z \in \mathfrak{r}\) ) be an element of the normalizer \(\mathfrak{u}\) of \(\mathfrak{h}\) in \(\mathfrak{g}\) . Show that \([y, \mathfrak{k}] \subset \mathfrak{k}\) , so \(y \in \mathfrak{k}\) and \(z \in \mathfrak{u}\) . Then show that \([z, \mathfrak{k}] \subset \mathfrak{h} \cap \mathfrak{r} \subset \mathfrak{c}\) , so \([\mathfrak{k}, [\mathfrak{k}, z]] = 0\) , and hence that \([\mathfrak{k}, z] = 0\) and \(z \in \mathfrak{c}\) . Finally, \([z, \mathfrak{l}] \subset \mathfrak{l}\) hence \(z \in \mathfrak{l}\) and \(x \in \mathfrak{h}\) .)
\item
  Let \(\mathfrak{s},\mathfrak{k},\mathfrak{c}\) be as in \(d\) ), and \(\mathfrak{q}=[\mathfrak{g},\mathfrak{r}]\) the nilpotent radical of \(\mathfrak{g}\) . Let \(x\in\mathfrak{q}\) and \(u\) the special automorphism \(e^{\mathrm{ad}x}\) . If \(u(\mathfrak{k})\subset\mathfrak{k}+\mathfrak{c}\) , then \(x\in\mathfrak{c}\) . (Consider the adjoint representation \(\rho\) of \(\mathfrak{s}\) on \(\mathfrak{q}\) , and let \(\mathfrak{q}_{i}\) be a complement of \(\mathcal{C}^{i+1}\mathfrak{q}\) in \(\mathcal{C}^{i}\mathfrak{q}\) stable under \(\rho\) ; let \(\rho_{i}\) be the subrepresentation of \(\rho\) defined by \(\mathfrak{q}_{i}\) . Let \(\sigma_{i}=\rho_{i}|\mathfrak{k}\) . Let \(\mathfrak{q}_{i}^{\prime}\) be the commutant of \(\mathfrak{k}\) in \(\mathfrak{q}_{i}\) and \(\mathfrak{q}_{i}^{\prime\prime}\) a complement of \(\mathfrak{q}_{i}^{\prime}\) in \(\mathfrak{q}_{i}\) stable under \(\sigma_{i}\) . Let \(x=x_{1}^{\prime}+x_{1}^{\prime\prime}+\cdots+x_{n}^{\prime}+x_{n}^{\prime\prime}\) with \(x_{i}^{\prime}\in\mathfrak{q}_{i}^{\prime},x_{i}^{\prime\prime}\in\mathfrak{q}_{i}^{\prime\prime}\) . Arguing by contradiction, assume that the \(x_{i}^{\prime\prime}\) are not all zero and \(x_{1}^{\prime\prime}=\cdots=x_{p-1}^{\prime\prime}=0\) , \(x_{p}^{\prime\prime}\neq0\) , for example. If \(h\in\mathfrak{k}\) , \(u(h)=h+[x_{p}^{\prime\prime},h]+y\) with \(y\in\mathcal{C}^{p+1}\mathfrak{q}\) . Since \(u(h)\in\mathfrak{k}+\mathfrak{c}\) , this gives \([x_{p}^{\prime\prime},h]+y\in\mathfrak{q}_{p}^{\prime}+\mathfrak{q}_{p+1}^{\prime}+\cdots+\mathfrak{q}_{n}^{\prime}\) , hence \([h,x_{p}^{\prime\prime}]\in\mathfrak{q}_{p}^{\prime}\) , so \([h,x_{p}^{\prime\prime}]=0\) . Then \(x_{p}^{\prime\prime}=0\) , a contradiction.)
\end{enumerate}

\(f\) ) Let \(\mathfrak{h}\) be a Cartan subalgebra of \(\mathfrak{g}\) . Then \(\mathfrak{h}\) can be expressed uniquely as the sum of \(\mathfrak{h} \cap \mathfrak{r}\) and a Cartan subalgebra of a Levi subalgebra of \(\mathfrak{g}\) . (For the uniqueness, use \(e\) ) and Th. 5 of Chap. I, §6, no. 8.) The Levi subalgebra in question is not unique in general.

\(g\) ) Let \(\mathfrak{h}\) be a Cartan subalgebra of \(\mathfrak{g}\) , and \(\mathfrak{t} = \mathfrak{g}^0 (\mathfrak{h} \cap \mathfrak{r})\) . Then \(\mathfrak{h}\) is a Cartan subalgebra of \(\mathfrak{t}\) . We have \(\mathfrak{g} = \mathfrak{t} + \mathfrak{r}\) (use a Fitting decomposition for the adjoint representation of \(\mathfrak{h} \cap \mathfrak{r}\) on \(\mathfrak{g}\) ). The algebra \(\mathfrak{t} \cap \mathfrak{r}\) is the radical of \(\mathfrak{t}\) and is nilpotent (use Exerc. 5 of §2).

\begin{enumerate}
\def\labelenumi{\alph{enumi})}
\setcounter{enumi}{7}
\tightlist
\item
  Assume that k is algebraically closed. Let h be a Cartan subalgebra of g. There exists a Levi subalgebra s of g such that, for all \(\lambda \in h^{*}\) ,
\end{enumerate}

\[
\mathfrak {g} ^ {\lambda} (\mathfrak {h}) = \left(\mathfrak {g} ^ {\lambda} (\mathfrak {h}) \cap \mathfrak {s}\right) + \left(\mathfrak {g} ^ {\lambda} (\mathfrak {h}) \cap \mathfrak {r}\right).
\]

(With the notations of g), take for s a Levi subalgebra of t such that \(\mathfrak{h} = (\mathfrak{h} \cap \mathfrak{s}) + (\mathfrak{h} \cap \mathfrak{r})\) ; this exists by b).) \(^{7}\)

¶ 12) a) Let g be a solvable Lie algebra, and G a finite subgroup (resp. compact subgroup if k = R or C) of Aut(g). Show that there exists a Cartan subalgebra of g stable under G. (Argue by induction on dim g, and reduce to the case in which g is an extension of a nilpotent algebra g/n by a commutative ideal n which is a non-trivial simple g/n-module (cf.~proof of Th. 3). The Cartan subalgebras of g then form an affine space attached to n, on which G operates. Conclude by an argument with the barycentre.)

\begin{enumerate}
\def\labelenumi{\alph{enumi})}
\setcounter{enumi}{1}
\tightlist
\item
  Let g be a solvable Lie algebra and s a Lie subalgebra of \(\text{Der}(\mathfrak{g})\) . Assume that the s-module g is semi-simple. Show that there exists a Cartan subalgebra of g stable under s. (Same method.)
\end{enumerate}

¶ 13) Let g be a Lie algebra and G a finite subgroup of Aut(g). Assume that G is hyper-solvable (Algebra, Chap. I, §6, Exerc. 26). Show that there exists a Cartan subalgebra of g stable under G.

(Argue by induction on dim g. Using Exerc. 12, reduce to the case where g is semi-simple. If G ≠ \{1\}, choose a normal subgroup C of G that is cyclic of prime order (Algebra, loc. cit.). The subalgebra s consisting of the elements invariant under C is reductive in g (§1, no. 5), and distinct from g. By the induction hypothesis, s has a Cartan subalgebra a stable under G. We have s ≠ 0 (Chap. I, §4, Exerc. 21 c)), so a ≠ 0. The commutant z of a in g is distinct from g and stable under G; choose a Cartan subalgebra h of z stable under G, and show that h is a Cartan subalgebra of g, cf.~the Cor. to Prop. 3.)

Construct a finite group of automorphisms of \(\mathfrak{sl}(2,\mathbf{C})\) which is isomorphic to \(A_{4}\) (and hence solvable) and which does not leave any Cartan subalgebra stable.

\begin{enumerate}
\def\labelenumi{\arabic{enumi})}
\setcounter{enumi}{13}
\item
  *Show that every irreducible complex (resp. real) linear representation of a hyper-solvable finite group G is induced by a representation of degree 1 (resp. of degree 1 or 2) of a subgroup of G. (Apply Exerc. 13 to the Lie algebras \(\mathfrak{gl}(n,\mathbf{C})\) and \(\mathfrak{gl}(n,\mathbf{R}).\) )
\item
  Assume that k is algebraically closed. Let g be a Lie algebra, h a Cartan subalgebra of g, and A a subset of g. Assume that A is dense in g (in the Zariski topology) and stable under \(\operatorname{Aut}_{e}(\mathfrak{g})\) . Show that \(A \cap h\) is dense in h. (Let X be the closure of \(A \cap h\) , and \(U = h - X\) . Assume that \(U \neq \varnothing\) . With the notations in Lemma 2, the image under F of \(U \times \mathfrak{g}^{\lambda_{1}}(\mathfrak{h}) \times \cdots \times \mathfrak{g}^{\lambda_{p}}(\mathfrak{h})\) contains a non-empty open subset of g. Since this image is contained in g-A, this contradicts the fact that A is dense in g.)
\item
  Let V be a finite dimensional k-vector space, and g a Lie subalgebra of \(\mathfrak{gl}(V)\) . We are going to show that the following three properties are equivalent: (i) The Cartan subalgebras of g are commutative and consist of semi-simple elements.
\end{enumerate}

\begin{enumerate}
\def\labelenumi{(\roman{enumi})}
\setcounter{enumi}{1}
\item
  Every regular element of \(\mathfrak{g}\) is semi-simple.
\item
  The semi-simple elements of g are dense in g in the Zariski topology.
\end{enumerate}

\begin{enumerate}
\def\labelenumi{\alph{enumi})}
\item
  Show that (i) \(\Longrightarrow\) (ii) \(\Longrightarrow\) (iii).
\item
  Let A be the set of semi-simple elements of \(\mathfrak{g}\) . Show that A is stable under \(\mathrm{Aut}_e(\mathfrak{g})\) .
\item
  Show that (iii) \(\Longrightarrow\) (i). (Reduce (App. I, Exerc. 1) to the case where k is algebraically closed. If h is a Cartan subalgebra of g, show, by using Exerc. 15, that \(A \cap h\) is dense in h; since \([x, y] = 0\) if \(x \in A \cap h, y \in h\) , deduce that h is commutative, from which (i) is immediate.)
\item
  Assume that k is R, C or a non-discrete complete ultrametric field of characteristic zero. Give g the topology defined by that of k. Show that properties (i), (ii), (iii) are equivalent to the following:
\end{enumerate}

\begin{enumerate}
\def\labelenumi{(\roman{enumi})}
\setcounter{enumi}{3}
\tightlist
\item
  The semi-simple elements of \(\mathfrak{g}\) are dense in \(\mathfrak{g}\) .
\end{enumerate}

(Show that (iv) \(\Longrightarrow\) (iii) and (ii) \(\Longrightarrow\) (iv), cf.~App. I, Exerc. 4.)

\begin{enumerate}
\def\labelenumi{\arabic{enumi})}
\setcounter{enumi}{16}
\tightlist
\item
  Assume that \(k\) is algebraically closed. Let \(\mathfrak{g}\) be a Lie algebra, \(\mathfrak{h}\) a Cartan subalgebra, and A a subset of the centre of \(\mathfrak{h}\) . Denote by \(\mathrm{E}_{\mathfrak{g}}\) the subgroup of \(\operatorname{Aut}(\mathfrak{g})\) denoted by E in no. 2. Show that, if \(s\) is an element of \(\operatorname{Aut}(\mathfrak{g})\) such that \(s\mathrm{A} = \mathrm{A}\) , there exists \(t \in \mathrm{E}_{\mathfrak{g}}\) such that \(t\mathfrak{h} = \mathfrak{h}\) and \(t|\mathrm{A} = \mathrm{Id}_{\mathrm{A}}\) ; in particular, \(ts|_{\mathrm{A}} = s|_{\mathrm{A}}\) . (Let \(\mathfrak{a}\) be the commutant of A in \(\mathfrak{g}\) ; since \(\mathfrak{h}\) and \(sh\) are Cartan subalgebras of \(\mathfrak{a}\) , there exists \(\theta \in \mathrm{E}_{\mathfrak{a}}\) such that \(\theta(s\mathfrak{h}) = \mathfrak{h}\) ; choose \(t\) from the elements of \(\mathrm{E}_{\mathfrak{g}}\) that extend \(\theta\) .)
\end{enumerate}

¶ 18) Let g be a Lie algebra, h a Cartan subalgebra of g and Ug (resp. Uh) the enveloping algebra of g (resp. h). A linear form \(\varphi\) on Ug is said to be central if it vanishes on {[}Ug, Ug{]}, i.e.~if \(\varphi(a.b) = \varphi(b.a)\) for all \(a, b \in Ug\) .

\(a)\) Let \(x, y \in \mathfrak{g}\) . Assume that there exists \(s \in \operatorname{Aut}_e(\mathfrak{g})\) such that \(s(x) = y\) . Show that

\[
\varphi (x ^ {n}) = \varphi (y ^ {n})
\]

for all \(n \in \mathbf{N}\) and for every central linear form \(\varphi\) on \(\mathrm{Ug}\) .

\begin{enumerate}
\def\labelenumi{\alph{enumi})}
\setcounter{enumi}{1}
\item
  Let \(\varphi\) be a central linear form on Ug whose restriction to Uh vanishes. Show that \(\varphi = 0\) . (We can assume that k is algebraically closed. Deduce from a) that we then have \(\varphi(x^{n}) = 0\) for all \(n \in N\) and all regular \(x \in g\) ; use a density argument to remove the assumption of regularity.)
\item
  Show that \(Ug = [Ug, Ug] + Uh\) .
\end{enumerate}

\(d\) ) Let V be a semi-simple \(\mathfrak{g}\) -module. Show that V is semi-simple as an \(\mathfrak{h}\) -module. In particular, \(V^{\lambda}(\mathfrak{h}) = V_{\lambda}(\mathfrak{h})\) for all \(\lambda \in \mathfrak{h}^{*}\) .

\begin{enumerate}
\def\labelenumi{\alph{enumi})}
\setcounter{enumi}{4}
\tightlist
\item
  Let \(V'\) be a semi-simple g-module. Assume that V and \(V'\) are isomorphic as h-modules. Show that they are isomorphic as g-modules. (If \(a \in Uh\) , remark that \(\operatorname{Tr}(a_{\mathrm{V}}) = \operatorname{Tr}(a_{\mathrm{V}'})\) ). Deduce, by using b) or c), that \(\operatorname{Tr}(x_{\mathrm{V}}) = \operatorname{Tr}(x_{\mathrm{V}'})\) for all \(x \in Ug\) , and conclude by using Algebra, Chap. VIII.)
\end{enumerate}

If k is algebraically closed, the assumption that ``V and V' are h-isomorphic'' is equivalent to saying that \(\dim\mathrm{V}_{\lambda}(\mathfrak{h})=\dim\mathrm{V}_{\lambda}^{\prime}(\mathfrak{h})\) for all \(\lambda\in\mathfrak{h}^{*}\) .

